\chapter*{Conclusion générale}
\addcontentsline{toc}{chapter}{Conclusion générale}
\markboth{Conclusion générale}{Conclusion générale}

Ce mémoire, réalisé dans le cadre du projet Tandem, portait sur la planification des trajectoires d'un robot parallèle à câbles destiné à la surveillance du stress hydrique des cultures. L'objectif était de définir un parcours de prise de vue qui couvre la parcelle, respecte les limites mécaniques du robot et réduit le temps de déplacement ainsi que l'énergie mécanique.

Après avoir posé le cadre théorique des robots à câbles et des méthodes d'optimisation multi-objectif, le travail a appliqué l'algorithme NSGA-II au robot Tandem à 8 câbles, au moyen d'une chaîne de traitement en quatre blocs implémentée en Python. Le dimensionnement de l'algorithme, fondé sur 50 exécutions et 12 exécutions de robustesse, a conduit à retenir 500 individus et 2000 générations pour la grille $7\times7$, et 600 individus et 5000 générations pour la grille $9\times9$.

Sur la zone de service sécurisée, tous les points sont atteignables et les tensions restent inférieures à la borne de 58~N (maximum de 35,7~N). Par rapport à un parcours en zigzag, les trajectoires optimisées réduisent le temps de déplacement de 22,7~\% (grille $7\times7$) et de 21,0~\% (grille $9\times9$), et l'énergie mécanique de 27,6~\% et de 28,7~\%. L'extension de la zone de travail est en revanche pénalisante : elle augmente l'énergie de 28~\% et de 116~\% selon la grille, et fait perdre quatre points en $9\times9$. La zone sécurisée est donc retenue.

Ces résultats restent issus de simulations, avec des câbles rigides, sans vent et avec une énergie quasi-statique. La robustesse n'a été évaluée que sur trois graines et les tensions n'ont pas été calculées sur les zones étendues. Les perspectives consistent à consolider l'étude numérique, à valider les trajectoires sur le prototype du SP1, puis à intégrer une replanification en temps réel. Les trajectoires obtenues constituent l'entrée du SP3, chargé de l'analyse des images.
